\chapter{Conclusiones y trabajo futuro}
\label{cap:10}

Este capítulo presenta las conclusiones generales derivadas del diseño, la implementación y la evaluación del sistema de apoyo a la decisión para la priorización de incidentes 112 en Castilla y León. En primer lugar, se analiza el grado de consecución de los objetivos generales y específicos propuestos al inicio del trabajo. A continuación, se detallan las principales contribuciones técnicas y metodológicas del proyecto, seguidas de una discusión sobre los resultados obtenidos, las limitaciones detectadas en esta fase de desarrollo y las líneas prioritarias de trabajo futuro para consolidar la arquitectura propuesta.

\section{Respuesta al objetivo del TFM}

El objetivo de este TFM era diseñar, implementar y evaluar un sistema de apoyo a la
decisión para la priorización inicial de incidentes 112 en Castilla y León, manteniendo
explicabilidad, trazabilidad, supervisión humana y prudencia metodológica. La respuesta
obtenida es positiva dentro del alcance de un prototipo académico: el sistema no
sustituye al operador ni pretende certificar prioridad operativa oficial, pero demuestra
que es viable construir una cadena reproducible que transforma registros reales en
recomendaciones P1--P4 revisables. Asimismo, este desarrollo aporta evidencia técnica, dentro de un entorno académico controlado, sobre la integración de modelos interpretables y explicaciones en lenguaje natural.

La aportación central no reside en una única métrica, sino en la coherencia del conjunto:
Registro 112 CyL como fuente empírica, limpieza trazable, variables V01--V15,
supervisión débil, control de fugas de información, Capa 2 interpretable, explicación
normativa local, contratos de datos y registro de la decisión humana.

Esta respuesta cierra los objetivos formulados en el Capítulo~\ref{cap:4}. El trabajo
analiza el dominio 112 y su marco normativo, define variables y prioridad académica,
audita el conjunto de datos de Castilla y León, construye etiquetas débiles, compara alternativas interpretables,
integra una capa explicativa y evalúa el resultado con métricas, errores críticos y
trazabilidad. Por tanto, la defensa del TFM no depende de una promesa de despliegue
operativo, sino de una cadena de evidencia verificable.

Además, el sistema se ha comprobado como prototipo ejecutable de
extremo a extremo: Capa 1 extrae señales desde texto, Capa 2 prioriza con RuleFit-lite,
Capa 3 explica con Ollama local y el backend devuelve una respuesta consumible por la
interfaz. Tras integrar la Capa 1, se ha regenerado además una prueba técnica
reproducible Capa 1--Capa 2--Capa 3 que supera 5/5 escenarios curados y confirma que el
flujo integrado no rompe los resultados de la Capa 2. Este hito aporta evidencia de compatibilidad entre capas dentro de los escenarios ensayados, sin equivaler a una validación operativa.

\section{Trazabilidad de objetivos específicos}

Para evidenciar el grado de consecución del trabajo, la
Tabla~\ref{tab:trazabilidad-oe} enlaza cada objetivo específico definido en el
Capítulo~\ref{cap:4} con su resultado y la evidencia que lo respalda. Todos los
objetivos planteados se consideran alcanzados dentro del alcance académico definido.

\begin{table}[htbp]
\centering
\scriptsize
\setlength{\tabcolsep}{4pt}
\begin{tabular}{|p{1.0cm}|p{1.4cm}|p{6.6cm}|p{3.4cm}|}
\hline
\textbf{OE} & \textbf{Estado} & \textbf{Grado de consecución} & \textbf{Evidencia} \\ \hline
OE1 & Logrado & Se analiza el problema de priorización temprana 112, necesidades del operador, restricciones de información y riesgos de automatización. & Caps.~\ref{cap:1}, \ref{cap:2} \\ \hline
OE2 & Logrado & Se revisa el marco normativo estatal y autonómico CyL y se construye la tabla de trazabilidad norma--criterio--variable--regla. & Cap.~\ref{cap:3} \\ \hline
OE3 & Logrado & Se define la taxonomía operativa, la escala académica P1--P4 y las variables V01--V15 con separación anti-\textit{leakage}. & Anexos~\ref{cap:anexoa}, \ref{cap:anexob}, \ref{cap:anexoc} \\ \hline
OE4 & Logrado & Se audita el Registro 112 CyL (9.380 registros limpios) y se generan etiquetas P1--P4 por supervisión débil con acuerdo Krippendorff \(\alpha = 0{,}899902\). & Cap.~\ref{cap:7}, Anexo~\ref{cap:anexoc} \\ \hline
OE5 & Logrado & Se implementa la Capa 2 interpretable: línea base heurística, RuleFit-lite y alternativa externa; se selecciona RuleFit-lite (macro-F1 0,905604; sensibilidad P1 0,909924). & Cap.~\ref{cap:9}, Anexos~\ref{cap:anexoe}, \ref{cap:anexol} \\ \hline
OE6 & Logrado & Se diseña la Capa 3 de explicación con LLM local, RAG normativo y MCP, sin delegar la decisión de prioridad en el LLM. & Cap.~\ref{cap:8}, Anexo~\ref{cap:anexol} \\ \hline
OE7 & Logrado & Se integran las capas mediante contratos de datos, trazabilidad de la inferencia, respuesta alternativa determinista y registro de la decisión humana. & Cap.~\ref{cap:8} \\ \hline
OE8 & Logrado & Se evalúa con particiones estratificada y temporal, control de fugas, 59 falsos negativos P1, fidelidad explicativa y un estudio exploratorio con nueve participantes. & Cap.~\ref{cap:9}, Anexo~\ref{cap:anexod} \\ \hline
\end{tabular}
\caption{Trazabilidad entre objetivos específicos, grado de consecución y evidencia.}
\label{tab:trazabilidad-oe}
\end{table}

\noindent\textit{Fuente: elaboración propia.}

\section{Principales contribuciones}

A lo largo de su desarrollo, este proyecto ha permitido consolidar una serie de aportes tanto teóricos como prácticos en el campo de los sistemas de soporte a la decisión en emergencias. Estas aportaciones abordan desde la estructuración de datos abiertos hasta la implementación de una arquitectura modular basada en la interpretabilidad y la supervisión humana. En particular, el trabajo realizado aporta siete contribuciones principales que se detallan a continuación.

\begin{enumerate}
    \item \textbf{Arquitectura en tres capas.} Se ha definido una arquitectura formada
    por Capa 1 NLP, Capa 2 RuleFit-lite y Capa 3 LLM+RAG+MCP, separando extracción,
    priorización y explicación normativa.
    \item \textbf{Uso realista de datos abiertos.} El sistema se apoya en 9.380
    registros limpios del Registro 112 de Castilla y León 2008--2022, con auditoría de
    calidad, cobertura temporal y control de campos ex post.
    \item \textbf{Etiqueta académica P1--P4 mediante supervisión débil.} Ante la
    ausencia de prioridad oficial en el conjunto de datos, se construye una etiqueta académica
    con cuatro funciones heurísticas correlacionadas. El acuerdo alcanza Krippendorff
    \(\alpha = 0{,}899902\), y la ablación sin reglas heurísticas mantiene
    \(\alpha = 0{,}834155\).
    \item \textbf{Modelo interpretable seleccionado.} RuleFit-lite supera a la línea base
    heurística en macro-F1, mantiene la sensibilidad P1 por encima del criterio fijado y
    conserva 30 reglas activas exportadas.
    \item \textbf{Evaluación reproducible y anti-leakage.} Los resultados se obtienen
    sobre particiones reproducibles, con prueba estratificada y prueba temporal, evitando columnas
    posteriores a la decisión.
    \item \textbf{Prototipo integrado local.} Se ha comprobado una primera cadena
    ejecutable con API, interfaz, RuleFit-lite y Ollama local
    \texttt{llama3.1:8b-instruct-q4\_K\_M}, sin delegar la prioridad al LLM. También se
    ha comprobado la integración reproducible Capa 1--Capa 2--Capa 3 sobre cinco
    escenarios curados, con RuleFit activo y salida coherente en todos los casos.
    \item \textbf{Trazabilidad y evaluación complementaria.} Se
    documenta una revisión conjunta de los tres autores sobre 119 casos, incluyendo
    los 59 falsos negativos P1 detectados en la prueba estratificada, y se evalúa la
    fidelidad de las explicaciones de Capa 3 sobre esa misma muestra. Además, nueve
    participantes aportan una primera valoración descriptiva de la utilidad y la
    comprensibilidad del prototipo.
\end{enumerate}

\section{Resultados principales}

La evaluación confirma que RuleFit-lite es la alternativa más equilibrada para la Capa
2. En la partición de prueba estratificada alcanza exactitud 0.903648, macro-F1
0.905604 y sensibilidad P1 0.909924. Estos valores superan a la línea base heurística de
19 reglas en macro-F1, que obtiene 0.464284, aunque la línea base conserva una sensibilidad P1
de 0.998473 por su
comportamiento conservador.

RuleFit-lite se ajusta con un submuestreo estratificado y reproducible de 3.000 casos de
los 6.512 disponibles. La ejecución de \texttt{imodels} usa 50 casos y se conserva como
prueba diagnóstica, de modo que no permite concluir que RuleFit-lite sea más eficiente
en condiciones equivalentes.

Estas métricas corresponden a RuleFit-lite sin la Puerta de Seguridad. La salvaguarda
determinista pertenece a la inferencia integrada y se comprueba de forma separada
mediante escenarios funcionales, para no atribuir al estimador los efectos del
posprocesamiento.

La prueba temporal sobre registros de 2022 mantiene resultados estables: exactitud
0.884354, macro-F1 0.877594 y sensibilidad P1 0.928767. Esta estabilidad no demuestra validez
prospectiva en producción, pero reduce el riesgo de que el resultado dependa solo de una
partición estratificada favorable.

El análisis de errores identifica 59 falsos negativos P1: 34 clasificados como P2, 18
como P3 y 7 como P4. La presencia de 7 casos P1 infrapriorizados como P4 representa un riesgo crítico de seguridad que queda documentado para su revisión. Este resultado es importante:
el sistema cumple el criterio de sensibilidad P1, pero los errores P1--P4 constituyen
el principal foco de mejora antes de cualquier validación externa.

La evaluación de fidelidad de explicaciones aporta una comprobación complementaria. Sobre
los 119 casos de la muestra revisada internamente, el juez determinista sin conexión obtiene una tasa de cumplimiento
1.000000, fidelidad media 0.957367 y fidelidad mínima 0.933333. En los 59 falsos
negativos P1, la fidelidad media es 0.952494. Esto no equivale a una evaluación externa
de calidad lingüística, pero aporta evidencia de que las explicaciones deterministas
evaluadas no contradicen la recomendación de Capa 2 y conservan la trazabilidad definida.

El cierre documental de la evaluación se recoge en un reporte maestro de trazabilidad que actúa como
índice de métricas, evidencias, revisión interna, fidelidad
explicativa y capítulos/anexos que sostienen los resultados.

La evaluación exploratoria con nueve participantes aporta una perspectiva complementaria
sobre utilidad percibida y comprensibilidad. Las medias globales son 4,14 para la
razonabilidad de la prioridad, 4,04 para la claridad de la explicación, 4,04
para la confianza aportada por señales o reglas y 4,63 para la claridad de la
decisión humana. El resultado más sólido es, por tanto, la comprensión del control humano;
la principal oportunidad de mejora es simplificar la presentación de probabilidades y
reglas para perfiles no especializados. Por el tamaño y la selección de la muestra, estas
puntuaciones son descriptivas y no acreditan aceptación operativa.

Como comprobación integrada del prototipo, la prueba funcional P1--P4 obtiene 4 aciertos en
4 casos representativos: incendio con atrapados como P1, fuga de gas sin heridos como
P2, árbol en calzada como P3 y consulta vecinal sin riesgo como P4. La prueba confirma
RuleFit activo, el LLM local disponible y explicación generativa. La latencia media
observada en el entorno experimental es \(4{,}898\) s, valor que caracteriza la integración funcional, pero queda
pendiente de optimización si se quisiera una demo más fluida.

La prueba reproducible Capa 1--Capa 2--Capa 3 complementa esta comprobación con 5
casos curados y resultado 5/5. Su valor principal es demostrar que la salida real de
Capa 1 alimenta correctamente a RuleFit-lite, incluyendo un caso ordinario con
negaciones explícitas que antes podía provocar falsos positivos de atrapamiento. Las
latencias medias son 4.58844 ms en Capa 1, 52.16898 ms en Capa 2 y 0.30748 ms en Capa
3 degradada.

\section{Cumplimiento de principios de diseño}

El prototipo respeta la línea de trabajo definida desde los capítulos iniciales. La
prioridad no la decide un LLM generativo, sino una Capa 2 interpretable. El LLM se
reserva para explicar y citar normativa. La inferencia se concibe en local, sin enviar
incidentes a APIs externas. La arquitectura mantiene contratos versionados y registra
la decisión humana mediante \texttt{OperatorDecision} e \texttt{InferenceLog}.

También se cumple el principio anti-\textit{leakage}: se excluyen los campos relativos
a medios movilizados, pacientes atendidos, cierre del incidente y actualizaciones
posteriores. Esta decisión es esencial para que el modelo represente el momento de
primera decisión y no información generada durante el operativo.

La Capa 1 se mantiene deliberadamente acotada. Aplica NLP simbólico mediante detección
semántica basada en reglas y diccionarios operativos, tratamiento de negaciones,
normalización y estimación de confianza para transformar texto libre en señales
auditables. Su evaluación oficial se limita a cinco escenarios curados y, por
tanto, acredita integración funcional, no rendimiento generalizable. Aunque el código
contempla una vía experimental para cargar un clasificador aprendido, no se dispone de
un modelo congelado ni de una evaluación suficiente para presentarlo como componente
validado. Esta delimitación preserva la coherencia entre memoria, código y evidencia.

\section{Limitaciones}

A pesar de los resultados favorables obtenidos en la fase experimental, el prototipo actual presenta restricciones que impiden interpretar sus resultados como evidencia de viabilidad inmediata en producción. Documentar estas limitaciones aporta transparencia y orienta las futuras iteraciones del sistema:

\begin{enumerate}
    \item La etiqueta P1--P4 es académica y débil; no existe una prioridad oficial
    equivalente en el Registro 112 CyL.
    \item La validación técnica es interna. La evaluación exploratoria incorpora nueve
    participantes, pero no un panel externo de operadores o responsables del 112.
    \item RuleFit-lite se ha ajustado con 3.000 de los 6.512 casos disponibles y la
    ejecución \texttt{imodels} con 50. Esta última se limita a una configuración
    diagnóstica y no sustenta una comparación de eficiencia en condiciones equivalentes.
    \item La calibración ECE queda pendiente de una pasada final si el modelo se
    congela para una demostración cerrada.
    \item Las variables V08--V11, asociadas a AEMET histórico, inundabilidad SNCZI y
    proximidad Seveso, permanecen diferidas como ampliación futura.
    \item El prototipo no mide impacto operativo real sobre tiempos de decisión,
    coordinación de recursos o carga cognitiva en sala.
    \item Las pruebas integradas tienen tamaño reducido: cuatro casos con Ollama y
    cinco escenarios reproducibles Capa 1--Capa 2--Capa 3. La evaluación exploratoria
    añade nueve participantes seleccionados intencionalmente, pero no incluye operadores
    del 112 ni equivale a una validación de aceptación operativa.
    \item La muestra de revisión interna contiene 119 casos, incluidos los 59 falsos
    negativos P1, y ha sido examinada conjuntamente por los tres autores. Al no existir
    tres anotaciones independientes, no permite calcular acuerdo interrevisor ni sustituye
    una validación externa.
\end{enumerate}

\section{Trabajo futuro}

La continuación natural del proyecto se organiza en siete líneas.

\begin{enumerate}
    \item \textbf{Validación externa con personal del 112.} Revisar una muestra de
    casos con operadores o responsables expertos, comparando etiqueta académica,
    recomendación del sistema y criterio profesional.
    \item \textbf{Revisión independiente.} Registrar por separado el criterio de cada
    revisor sobre los 119 casos del Anexo~\ref{cap:anexod}, calcular acuerdo entre autores
    y documentar patrones de divergencia.
    \item \textbf{Calibración probabilística.} Evaluar ECE y aplicar calibración si el
    modelo se congela como demostrador estable.
    \item \textbf{Variables contextuales avanzadas.} Incorporar AEMET histórico,
    SNCZI/zonas inundables y Seveso como variables V08--V11, manteniendo el control
    anti-\textit{leakage}.
    \item \textbf{Evaluación de explicaciones con juez independiente.} Desarrollar la
    propuesta de implementar un LLM-as-Judge local independiente para evaluar de forma
    ciega las explicaciones de la Capa 3, usando la muestra de 119 casos
    y comparando sus puntuaciones de fidelidad y calidad lingüística con el juez
    determinista offline.
    \item \textbf{Optimización del prototipo integrado y entrada de voz.} Reducir la latencia del LLM local mediante instrucciones más compactas, respuesta progresiva, cuantización o ejecución acelerada; evaluar experimentalmente el reconocimiento automático del habla con \texttt{faster-whisper} sobre un corpus autorizado y representativo antes de atribuirle validez sobre llamadas reales.
    \item \textbf{Piloto controlado ampliado.} Extender la evaluación del prototipo en
    un entorno de prueba controlado y con una muestra representativa para medir utilidad
    percibida, carga cognitiva, tiempo de revisión y aceptación de recomendaciones.
\end{enumerate}

\section{Cierre}

El TFM alcanza una conclusión equilibrada: es posible construir un sistema inteligente
de apoyo a la priorización 112 que sea técnicamente ambicioso y, al mismo tiempo,
prudente en un dominio crítico. RuleFit-lite aporta rendimiento superior a la línea base
experto sin renunciar a interpretabilidad; la prueba temporal aporta una primera evidencia
de estabilidad; y la revisión interna conjunta abre el camino para contrastar el sistema con
criterio humano.

El prototipo no es un producto operativo, pero constituye una base académica sólida de investigación
aplicada. Su valor está en haber unido datos reales, normativa, aprendizaje
interpretable, explicación local y supervisión humana en una arquitectura reproducible y
auditable. Esa combinación permite defender el proyecto por su rigor metodológico y por
delimitar con claridad hasta dónde puede llegar la IA responsable en una primera
decisión de emergencia.
